\documentclass[fontset=windows]{ctexart}

\usepackage{anyfontsize}
\usepackage{amsmath,amssymb,mathrsfs,bm}

\allowdisplaybreaks[4]
\usepackage{indentfirst}
\setlength{\parindent}{0em}
\usepackage{geometry}
\geometry{top=3cm,bottom=3cm,left=2cm,right=2cm}

\begin{document}

    \title{\textbf{数学物理方法作业}}
    \author{1900011604\ \ 张植竣\ \ 北京大学}
    \date{2020年11月10日}

    \maketitle

    \section*{第十六章}

    \bigskip
    \begin{itemize}

    \item[2.]
    代入Legendre多项式的微分表示，分部积分$l$次，得：
    \begin{align*}
            \int_{-1}^{1}(1+x)^k\mathrm{P}_l(x)\,\mathrm{d}x
            &= \frac{1}{2^l l!}\int_{-1}^{1}\frac{\mathrm{d}(1+x)^k}{\mathrm{d}x}\frac{\mathrm{d}^l(x^2-1)^l}{\mathrm{d}x^l}\,\mathrm{d}x \\
            &= (-1)^l\frac{1}{2^l l!}\int_{-1}^{1}\frac{\mathrm{d}^l(1+x)^k}{\mathrm{d}x^l}(x^2-1)^{l}\,\mathrm{d}x \\
            &= \frac{1}{2^l l!}\int_{-1}^{1}\frac{\mathrm{d}^l(1+x)^k}{\mathrm{d}x^l}(1-x^2)^{l}\,\mathrm{d}x
    \end{align*}
    分两种情况讨论：

    $k<l$时，应有：
    \[
        \frac{\mathrm{d}^l(1+x)^k}{\mathrm{d}x^l} = 0
    \]
    则：
    \[
        \int_{-1}^{1}(1+x)^k\mathrm{P}_l(x)\,\mathrm{d}x = 0
    \]
    $k\geqslant l$时，
    \begin{align*}
        \frac{1}{2^l l!}\int_{-1}^{1}\frac{\mathrm{d}^l(1+x)^k}{\mathrm{d}x^l}(1-x^2)^{l}\,\mathrm{d}x
        &= \frac{1}{2^l}\frac{k!}{l!(k-l)!}\int_{-1}^{1}(1+x)^{k-l}(1-x^2)^{l}\,\mathrm{d}x \\
        &= \frac{1}{2^l}\frac{k!}{l!(k-l)!}\int_{-1}^{1}(1+x)^{k-l}(1+x)^{l}(1-x)^{l}\,\mathrm{d}x \\
        &= \frac{1}{2^l}\frac{k!}{l!(k-l)!}\int_{-1}^{1}(1+x)^{k}(1-x)^{l}\,\mathrm{d}x
    \end{align*}
    再分部积分$l$次得：
    \begin{align*}
        \int_{-1}^{1}(1+x)^k(1-x)^{l}\,\mathrm{d}x
        &= \left.\frac{1}{k+1}(1+x)^{k+1}(1-x)^l\right|_{-1}^1 - \int_{-1}^{1}\frac{1}{k+1}(1+x)^{k+1}[-l(1-x)^{l-1}]\,\mathrm{d}x \\
        &= \left.\frac{1}{k+1}(1+x)^{k+1}(1-x)^l\right|_{-1}^1 + \frac{l}{k+1}\int_{-1}^{1}(1+x)^{k+1}(1-x)^{l-1}\,\mathrm{d}x \\
        &= \frac{l}{k+1}\int_{-1}^{1}(1+x)^{k+1}(1-x)^{l-1}\,\mathrm{d}x \\
        &= \frac{l}{k+1}\cdot\frac{l-1}{k+2}\int_{-1}^{1}(1+x)^{k+2}(1-x)^{l-2}\,\mathrm{d}x \\
        &= \frac{l!k!}{(k+l)!}\int_{-1}^{1}(1+x)^{k+l}\,\mathrm{d}x \\
        &= \left.\frac{l!k!}{(k+l)!}\cdot\frac{1}{k+l+1}(1+x)^{k+l+1}\right|_{-1}^{1} \\
        &= \frac{l!k!}{(k+l+1)!}2^{k+l+1}
    \end{align*}
    则原积分：
    \begin{align*}
        \int_{-1}^{1}(1+x)^k\mathrm{P}_l(x)\,\mathrm{d}x
        &= \frac{1}{2^l l!}\int_{-1}^{1}\frac{\mathrm{d}^l(1+x)^k}{\mathrm{d}x^l}(1-x^2)^{l}\,\mathrm{d}x \\
        &= \frac{1}{2^l }\frac{k!}{l!(k-l)!}\int_{-1}^{1}(1+x)^{k}(1-x)^{l}\,\mathrm{d}x \\
        &= \frac{1}{2^l }\frac{k!}{l!(k-l)!}\int_{-1}^{1}(1+x)^{k}(1-x)^{l}\,\mathrm{d}x \\
        &= \frac{1}{2^l}\frac{k!}{l!(k-l)!}\cdot\frac{l!k!}{(k+l+1)!}2^{k+l+1} \\
        &= \frac{(k!)^2}{(k-l)!(k+l+1)!}2^{k+1}
    \end{align*}
    综上：
    \[
    \int_{-1}^{1}(1+x)^k\mathrm{P}_l(x)\,\mathrm{d}x =
    \left\{
    \begin{aligned}
        &\ 0,\qquad &k<l \\
        &\ \frac{(k!)^2}{(k-l)!(k+l+1)!}2^{k+1},\qquad &k\geqslant l
    \end{aligned}
    \right.
    \]
    \bigskip

    \item[3.(1)]
    分两种情况讨论：

    $l=0$时，$\mathrm{P}_0(x) = 1$，直接积分：
    \begin{align*}
        \int_{-1}^{1}\mathrm{P}_l(x)\ln(1-x)\,\mathrm{d}x
        &= \int_{-1}^{1}\ln(1-x)\,\mathrm{d}x \\
        &= \left.(1-x)-(1-x)\ln(1-x)\right|_{-1}^{1} \\
        &= 2\ln2-2
    \end{align*}
    $l>0$时，代入Legendre方程，分部积分：
    \begin{align*}
        \int_{-1}^{1}\mathrm{P}_l(x)\ln(1-x)\,\mathrm{d}x
        &= \int_{-1}^{1}\ln(1-x)\cdot\left(-\frac{1}{l(l+1)}\right)\frac{\mathrm{d}}{\mathrm{d}x}\left[(1-x^2)\frac{\mathrm{d}\mathrm{P}_l(x)}{\mathrm{d}x}\right]\,\mathrm{d}x \\
        &= -\frac{1}{l(l+1)}\int_{-1}^{1}\ln(1-x)\frac{\mathrm{d}}{\mathrm{d}x}\left[(1-x^2)\frac{\mathrm{d}\mathrm{P}_l(x)}{\mathrm{d}x}\right]\,\mathrm{d}x \\
        &= -\frac{1}{l(l+1)}\int_{-1}^{1}\ln(1-x)\,\mathrm{d}\left[(1-x^2)\frac{\mathrm{d}\mathrm{P}_l(x)}{\mathrm{d}x}\right] \\
        &= \frac{1}{l(l+1)}\int_{-1}^{1}(1-x^2)\frac{\mathrm{d}\mathrm{P}_l(x)}{\mathrm{d}x}\mathrm{d}\ln(1-x)
    \end{align*}
    其中利用了：
    \begin{align*}
        \lim_{x\to1}\frac{1-x}{\ln(1-x)}
        &= \lim_{x\to1}\dfrac{-1}{-\dfrac{1}{1-x}} \\
        &= \lim_{x\to1}(1-x) \\
        &= 0
    \end{align*}
    则
    \[
    \left.\ln(1-x)\left[(1-x^2)\frac{\mathrm{d}\mathrm{P}_l(x)}{\mathrm{d}x}\right]\right|_{-1}^{1} = 0
    \]
    再次分部积分：
    \begin{align*}
        \int_{-1}^{1}\mathrm{P}_l(x)\ln(1-x)\,\mathrm{d}x
        &= \frac{1}{l(l+1)}\int_{-1}^{1}(1-x^2)\frac{\mathrm{d}\mathrm{P}_l(x)}{\mathrm{d}x}\mathrm{d}\ln(1-x) \\
        &= \frac{1}{l(l+1)}\int_{-1}^{1}(1-x^2)\frac{\mathrm{d}\mathrm{P}_l(x)}{\mathrm{d}x}\left(-\frac{1}{1-x}\,\mathrm{d}x\right) \\
        &= -\frac{1}{l(l+1)}\int_{-1}^{1}(1+x)\frac{\mathrm{d}\mathrm{P}_l(x)}{\mathrm{d}x}\,\mathrm{d}x \\
        &= -\frac{1}{l(l+1)}\int_{-1}^{1}(1+x)\,\mathrm{d}\mathrm{P}_l(x) \\
        &= -\left.\frac{1}{l(l+1)}(1+x)\mathrm{P}_l(x)\right|_{-1}^{1} + \frac{1}{l(l+1)}\int_{-1}^{1}\mathrm{P}_l(x)\,\mathrm{d}(1+x) \\
        &= -\left.\frac{1}{l(l+1)}(1+x)\mathrm{P}_l(x)\right|_{-1}^{1} + \frac{1}{l(l+1)}\int_{-1}^{1}\mathrm{P}_l(x)\,\mathrm{d}x \\
        &= -\frac{2}{l(l+1)}
    \end{align*}
    其中利用了：
    \[
    \mathrm{P}_l(1) = 1,\qquad \mathrm{P}_l(-1) = (-1)^l
    \]
    \[
    \int_{-1}^{1}x^n\mathrm{P}_l(x)\,\mathrm{d}x = 0,\qquad n<l
    \]
    综上：
    \[
    \int_{-1}^{1}\mathrm{P}_l(x)\ln(1-x)\,\mathrm{d}x =
    \left\{
    \begin{aligned}
        &\ 2\ln2-2,\qquad &l=0 \\
        &\ -\frac{2}{l(l+1)},\qquad &l>0
    \end{aligned}
    \right.
    \]
    \bigskip

    \item[3.(3)]
    对于Legendre多项式$\mathrm{P}_k(x)$，$k$为奇数时，$\mathrm{P}_k(x)$为奇函数，$k$为偶数时，$\mathrm{P}_k(x)$为偶函数。

    分两种情况讨论：

    $k+l=2a\ (a\in\mathbb{N})$时，$\mathrm{P}_k(x)\mathrm{P}_l(x)$为偶函数，直接积分：
    \begin{align*}
        \int_{0}^{1}\mathrm{P}_k(x)\mathrm{P}_l(x)\,\mathrm{d}x
        &= \frac{1}{2}\int_{-1}^{1}\mathrm{P}_k(x)\mathrm{P}_l(x)\,\mathrm{d}x \\
        &= \frac{1}{2}\cdot\frac{2}{2l+1}\delta_{kl} \\
        &= \frac{1}{2l+1}\delta_{kl}
    \end{align*}
    $k+l=2a+1\ (a\in\mathbb{N})$时，$\mathrm{P}_k(x)\mathrm{P}_l(x)$为奇函数。由Legendre方程得：
    \[
        \frac{\mathrm{d}}{\mathrm{d}x}\left[(1-x^2)\frac{\mathrm{d}\mathrm{P}_k(x)}{\mathrm{d}x}\right] + k(k+1)\mathrm{P}_k(x) = 0 \tag{1}
    \]
    \[
        \frac{\mathrm{d}}{\mathrm{d}x}\left[(1-x^2)\frac{\mathrm{d}\mathrm{P}_l(x)}{\mathrm{d}x}\right] + l(l+1)\mathrm{P}_l(x) = 0 \tag{2}
    \]
    $(1)\times \mathrm{P}_l(x) - (2)\times \mathrm{P}_k(x)$，并整理得：
    \[
    [k(k+1)-l(l+1)]\mathrm{P}_k(x)\mathrm{P}_l(x) = \mathrm{P}_k(x)\frac{\mathrm{d}}{\mathrm{d}x}\left[(1-x^2)\frac{\mathrm{d}\mathrm{P}_l(x)}{\mathrm{d}x}\right] - \mathrm{P}_l(x)\frac{\mathrm{d}}{\mathrm{d}x}\left[(1-x^2)\frac{\mathrm{d}\mathrm{P}_k(x)}{\mathrm{d}x}\right]
    \]
    上式两边对$x$积分得：
    \begin{align*}
        &\quad[k(k+1)-l(l+1)]\int_{0}^{1}\mathrm{P}_k(x)\mathrm{P}_l(x)\,\mathrm{d}x \\
        &= \int_{0}^{1}\left\{\mathrm{P}_k(x)\frac{\mathrm{d}}{\mathrm{d}x}\left[(1-x^2)\frac{\mathrm{d}\mathrm{P}_l(x)}{\mathrm{d}x}\right] - \mathrm{P}_l(x)\frac{\mathrm{d}}{\mathrm{d}x}\left[(1-x^2)\frac{\mathrm{d}\mathrm{P}_k(x)}{\mathrm{d}x}\right]\right\}\mathrm{d}x \\
        &= \left.(1-x^2)\left[\mathrm{P}_k(x)\frac{\mathrm{d}\mathrm{P}_l(x)}{\mathrm{d}x} - \mathrm{P}_l(x)\frac{\mathrm{d}\mathrm{P}_k(x)}{\mathrm{d}x}\right]\right|_{0}^{1} \\
        &= \mathrm{P}_l(0)\mathrm{P}'_k(0) - \mathrm{P}_k(0)\mathrm{P}'_l(0)
    \end{align*}
    由Legendre多项式：
    \[
        \mathrm{P}_n(x) = \sum_{r=0}^{[n/2]}(-1)^r\frac{(2n-2r)!}{2^n r!(n-r)!(n-2r)!}x^{n-2r}
    \]
    令$n=2m$，$r=m$得：
    \[
        \mathrm{P}_{2m}(0) = \frac{(-1)^m(2m)!}{(2^m m!)^2},\qquad m\in\mathbb{N}
    \]
    \[
        \mathrm{P}'_{2m}(0) = 0,\qquad m\in\mathbb{N}
    \]
    令$n=2m+1$，$r=m$得：
    \[
    \mathrm{P}_{2m+1}(0) = 0,\qquad m\in\mathbb{N}
    \]
    \[
    \qquad \mathrm{P}'_{2m+1}(0) = \frac{(-1)^m (2m+2)!}{2^{2m+1} m!(m+1)!},\qquad m\in\mathbb{N}
    \]
    由于$k+l=2a+1\ (a\in\mathbb{N})$，不妨设$k=2s$，$l=2t+1$，则：
    \begin{align*}
    \int_{0}^{1}\mathrm{P}_k(x)\mathrm{P}_l(x)\,\mathrm{d}x
    &= \frac{\mathrm{P}_l(0)\mathrm{P}'_k(0) - \mathrm{P}_k(0)\mathrm{P}'_l(0)}{2s(2s+1)-(2t+1)(2t+2)} \\
    &= \frac{1}{2s(2s+1)-(2t+1)(2t+2)}\left[0-\frac{(-1)^s(2s)!}{(2^s s!)^2}\cdot\frac{(-1)^t (2t+2)!}{2^{2t+1} t!(t+1)!}\right] \\
    &= \frac{1}{2s(2s+1)-(2t+1)(2t+2)}\cdot\frac{(-1)^{s+t+1}(2s)!(2t+2)!}{2^{2s+2t+1} (s!)^2 t!(t+1)!} \\
    &= \frac{-1}{2s(2s+1)-(2t+1)(2t+2)}\cdot\frac{(-1)^{s+t}(2s)!(2t+1)!(2t+2)}{2^{2s+2t} (s!t!)^2\cdot2(t+1)} \\
    &= \frac{(-1)^{s+t}}{(2t+1)(2t+2)-2s(2s+1)}\cdot\frac{(2s)!(2t+1)!}{2^{2s+2t} (s!t!)^2}
    \end{align*}
    综上：
    \[
    \int_{0}^{1}\mathrm{P}_k(x)\mathrm{P}_l(x)\,\mathrm{d}x =
    \left\{
    \begin{aligned}
        &\ \frac{(-1)^{s+t}\gamma_{kl}}{(2t+1)(2t+2)-2s(2s+1)}\cdot\frac{(2s)!(2t+1)!}{2^{2s+2t} (s!t!)^2},\qquad &k+l=2a+1 \\
        &\ \frac{1}{2l+1}\delta_{kl},\qquad &k+l=2a
    \end{aligned}
    \right.
    \]
    其中$a,s,t\in\mathbb{N}$，$s=\left[\dfrac{k}{2}\right]$，$t=\left[\dfrac{l}{2}\right]$，且
    \[
    \gamma_{kl} =
    \left\{
    \begin{aligned}
        &\ 1,\qquad &k=2s \\
        &\ -1,\qquad &k=2s+1
    \end{aligned}
    \right.
    \]
    \bigskip

    \item[3.(5)]
    由Legendre多项式的递推关系：
    \[
    (2l+1)x\mathrm{P}_l(x) = (l+1)\mathrm{P}_{l+1}(x)+l\mathrm{P}_{l-1}(x)
    \]
    得到：
    \[
    x\mathrm{P}_l(x) = \frac{(l+1)\mathrm{P}_{l+1}(x)+l\mathrm{P}_{l-1}(x)}{2l+1}
    \]
    \[
    x\mathrm{P}_{l+2}(x) = \frac{(l+3)\mathrm{P}_{l+3}(x)+(l+2)\mathrm{P}_{l+1}(x)}{2l+5}
    \]
    代入积分式得：
    \begin{align*}
        \int_{-1}^{1}x^2\mathrm{P}_l(x)\mathrm{P}_{l+2}(x)\,\mathrm{d}x
        &= \int_{-1}^{1}\frac{(l+1)\mathrm{P}_{l+1}(x)+l\mathrm{P}_{l-1}(x)}{2l+1}\cdot\frac{(l+3)\mathrm{P}_{l+3}(x)+(l+2)\mathrm{P}_{l+1}(x)}{2l+5}\,\mathrm{d}x \\
        &= \int_{-1}^{1}\frac{l+1}{2l+1}\cdot\frac{l+2}{2l+5}\mathrm{P}_{l+1}^2(x)\,\mathrm{d}x \\
        &= \frac{l+1}{2l+1}\cdot\frac{l+2}{2l+5}\cdot\frac{2}{2l+3} \\
        &= \frac{2(l+1)(l+2)}{(2l+1)(2l+3)(2l+5)}
    \end{align*}
    \bigskip

    \item[4.(1)]
    由Legendre多项式：
    \[
    \mathrm{P}_n(x) = \sum_{r=0}^{[n/2]}(-1)^r\frac{(2n-2r)!}{2^n r!(n-r)!(n-2r)!}x^{n-2r}
    \]
    可以得到：$\mathrm{P}_n(x)$为$n$阶多项式；且$n$为奇数时，$\mathrm{P}_n$为奇函数；$n$为偶数时，$\mathrm{P}_n$为偶函数。

    又因$f(x)=x^2$为偶函数，则其展开式应只含：
    \[
    \mathrm{P}_0(x) = 1,\qquad \mathrm{P}_2(x) = \frac{3}{2}x^2 - \frac{1}{2}
    \]
    即：
    \[
    f(x) = A\mathrm{P}_0(x) + B\mathrm{P}_2(x)
    \]
    解得：
    \[
    A = \frac{1}{3},\qquad B = \frac{2}{3}
    \]
    故：
    \[
    f(x) = x^2 = \frac{1}{3}\mathrm{P}_0(x) + \frac{2}{3}\mathrm{P}_2(x)
    \]
    \bigskip

    \item[4.(3)]
    考虑函数$f(x)=\mathrm{P}_l(x)\eta(x)$的Legendre多项式展开：
    \[
    f(x) = \sum_{k=0}^{\infty}c_k\mathrm{P}_k(x)
    \]
    其中：
    \begin{align*}
    c_k
    &= \frac{2k+1}{2}\int_{-1}^{1}f(x)\mathrm{P}_k(x)\,\mathrm{d}x \\
    &= \frac{2k+1}{2}\int_{0}^{1}\mathrm{P}_l(x)\mathrm{P}_k(x)\,\mathrm{d}x \\
    \end{align*}
    代入：
    \[
    \int_{0}^{1}\mathrm{P}_k(x)\mathrm{P}_l(x)\,\mathrm{d}x =
    \left\{
    \begin{aligned}
        &\ \frac{(-1)^{s+t}\gamma_{kl}}{(2t+1)(2t+2)-2s(2s+1)}\cdot\frac{(2s)!(2t+1)!}{2^{2s+2t} (s!t!)^2},\qquad &k+l=2a+1 \\
        &\ \frac{1}{2l+1}\delta_{kl},\qquad &k+l=2a
    \end{aligned}
    \right.
    \]
    其中$a,s,t\in\mathbb{N}$，$s=\left[\dfrac{k}{2}\right]$，$t=\left[\dfrac{l}{2}\right]$，且
    \[
    \gamma_{kl} =
    \left\{
    \begin{aligned}
        &\ 1,\qquad &k=2s \\
        &\ -1,\qquad &k=2s+1
    \end{aligned}
    \right.
    \]
    则：
    \[
    f(x) = \frac{1}{2}\mathrm{P}_l(x) + \sum_{k=0}^{\infty}\frac{2k
        +1}{2}\cdot\frac{(-1)^{s+t}\gamma_{kl}}{(2t+1)(2t+2)-2s(2s+1)}\cdot\frac{(2s)!(2t+1)!}{2^{2s+2t} (s!t!)^2}\mathrm{P}_k(x)
    \]
    令$l=1$，则$k$一定取偶数值，$t=0$，$\gamma_{kl}=1$，上式化为：
    \begin{align*}
        f(x)
        &= x\eta(x) \\
        &= \frac{x+|x|}{2} \\
        &= \frac{1}{2}x + \sum_{s=0}^{\infty}\frac{4s+1}{2}\cdot\frac{(-1)^{s}}{2-2s(2s+1)}\cdot\frac{(2s)!}{2^{2s} (s!)^2}\mathrm{P}_{2s}(x) \\
    \end{align*}
    得到：
    \[
    |x| = \sum_{s=0}^{\infty}(4s+1)\cdot\frac{(-1)^{s}}{2-2s(2s+1)}\cdot\frac{(2s)!}{2^{2s} (s!)^2}\mathrm{P}_{2s}(x)
    \]
    整理得：
    \[
    f(x) = |x| = \sum_{s=0}^{\infty}\frac{(-1)^{s-1}(2s)!}{\left(2^{s} s!\right)^2}\cdot\frac{(4s+1)}{4s^2+2s-2}\mathrm{P}_{2s}(x)
    \]
    \bigskip

    \item[5.]
    由题意，$u$只与$r$、$\theta$有关，与$\phi$无关，设$u=u(r,\theta)$。

    设方程的解具有分离变量的形式：
    \[
    u(r,\theta) = R(r)\varTheta(\theta)
    \]
    代入方程$\nabla^2 u=0$并分离变量得：
    \[
    \frac{\mathrm{d}}{\mathrm{d}r}\left(r^2\frac{\mathrm{d}R}{\mathrm{d}r}\right) - \lambda R = 0 \tag{1}
    \]
    \[
    \frac{1}{\sin\theta}\cdot\frac{\mathrm{d}}{\mathrm{d}\theta}\left(\sin\theta\frac{\mathrm{d}\varTheta}{\mathrm{d}\theta}\right) + \lambda\varTheta = 0 \tag{2}
    \]
    由$\theta$的边界条件，得：
    \[
    \varTheta(0)<\infty,\qquad \varTheta(\pi)<\infty \tag{3}
    \]
    (2)和(3)构成了Legendre本征值问题，解为：
    \[
    \lambda_l = l(l+1)
    \]
    \[
    \varTheta_l(\theta) = \mathrm{P}_l(\cos\theta)
    \]
    代入方程(1)得：
    \[
    \frac{\mathrm{d}}{\mathrm{d}r}\left(r^2\frac{\mathrm{d}R}{\mathrm{d}r}\right) - l(l+1)R = 0 \tag{4}
    \]
    整理得：
    \[
    r^2\frac{\mathrm{d}^2 R}{\mathrm{d}r^2} + 2r\frac{\mathrm{d}R}{\mathrm{d}r} - l(l+1)R = 0
    \]
    令$t = \ln(r)$，得：
    \[
    \left(\frac{\mathrm{d}^2 R}{\mathrm{d}t^2} - \frac{\mathrm{d}R}{\mathrm{d}t}\right) + 2\frac{\mathrm{d}R}{\mathrm{d}t} - l(l+1)R = 0
    \]
    特征方程为：
    \[
    \rho^2 + \rho - l(l+1) = 0
    \]
    解得：
    \[
    \rho_1 = l,\qquad \rho_2 = -l-1
    \]
    则：
    \[
    R_l(r) = A_lr^l + B_lr^{-l-1}
    \]
    定解问题的一般解为：
    \[
    u(r,\theta) = \sum_{l=0}^{\infty}(A_lr^l + B_lr^{-l-1})\mathrm{P}_l(\cos\theta)
    \]
    考虑边界条件：
    \[
    \left.u\right|_{r=a} = u_0,\qquad \left.u\right|_{r=b} = u_0\cos^2\theta = \frac{2}{3}u_0 \mathrm{P}_2(\cos\theta)+\frac{1}{3}u_0 \mathrm{P}_0(\cos\theta)
    \]
    得方程：
    \[
    \left\{
    \begin{aligned}
        & A_0 + B_0 a^{-1} = u_0 \\
        & A_2 a^2 + B_2 a^{-3} = 0 \\
        & A_0 + B_0 b^{-1} = \frac{1}{3}u_0 \\
        & A_2 b^2 + B_2 b^{-3} = \frac{2}{3}u_0
    \end{aligned}
    \right.
    \]
    解得：
    \[
    \left\{
    \begin{aligned}
        & A_0 = \frac{b-3a}{3(b-a)}u_0 \\
        & A_2 = \frac{2b^3 a^2}{3a^2(b^5-a^5)}u_0 \\
        & B_0 = \frac{2ab}{3(b-a)}u_0 \\
        & B_2 = \frac{2b^3 a^5}{3(b^5-a^5)}u_0
    \end{aligned}
    \right.
    \]
    且其余系数均为0。

    综上，定解问题的解为：
    \[
    u(r,\theta) = \frac{b-3a}{3(b-a)}u_0 + \frac{2b}{3(b-a)}\cdot\frac{a}{r}u_0 +\frac{2b^3 a^2}{3(b^5-a^5)}\left[\left(\frac{r}{a}\right)^2 - \left(\frac{a}{r}\right)^3\right]u_0 \mathrm{P}_2(\cos\theta)
    \]
    \bigskip

    \item[7.]
    设半球内的温度分布为$u = u(r,\theta,\phi)$，则可得定解问题：
    \[
    \nabla^2 u = 0,\qquad 0\leqslant\theta\leqslant\ \frac{\pi}{2}
    \]
    \[
    \left.u\right|_{r=0} < \infty,\qquad \left.u\right|_{r=a} = u_0\mathrm{sgn}(\cos\theta)
    \]
    由题意，$u$只与$r$、$\theta$有关，与$\phi$无关，设$u = u(r,\theta)$。

    设方程的解具有分离变量的形式：
    \[
    u(r,\theta) = R(r)\varTheta(\theta)
    \]
    代入方程$\nabla^2 u=0$并分离变量得：
    \[
    \frac{\mathrm{d}}{\mathrm{d}r}\left(r^2\frac{\mathrm{d}R}{\mathrm{d}r}\right) - \lambda R = 0 \tag{1}
    \]
    \[
    \frac{1}{\sin\theta}\cdot\frac{\mathrm{d}}{\mathrm{d}\theta}\left(\sin\theta\frac{\mathrm{d}\varTheta}{\mathrm{d}\theta}\right) + \lambda\varTheta = 0 \tag{2}
    \]
    由$\theta$的边界条件，得：
    \[
    \varTheta(0)<\infty,\qquad \varTheta\left(\frac{\pi}{2}\right)=0 \tag{3}
    \]
    (2)和(3)构成了Legendre本征值问题，其解为：
    \[
    \lambda_l = l(l+1)
    \]
    \[
    \varTheta_l(\theta) = \mathrm{P}_l(\cos\theta)
    \]
    代入方程(1)得：
    \[
    \frac{\mathrm{d}}{\mathrm{d}r}\left(r^2\frac{\mathrm{d}R}{\mathrm{d}r}\right) - l(l+1)R = 0 \tag{4}
    \]
    整理得：
    \[
    r^2\frac{\mathrm{d}^2 R}{\mathrm{d}r^2} + 2r\frac{\mathrm{d}R}{\mathrm{d}r} - l(l+1)R = 0
    \]
    令$t = \ln(r)$，得：
    \[
    \left(\frac{\mathrm{d}^2 R}{\mathrm{d}t^2} - \frac{\mathrm{d}R}{\mathrm{d}t}\right) + 2\frac{\mathrm{d}R}{\mathrm{d}t} - l(l+1)R = 0
    \]
    特征方程为：
    \[
    \rho^2 + \rho - l(l+1) = 0
    \]
    解得：
    \[
    \rho_1 = l,\qquad \rho_2 = -l-1
    \]
    则：
    \[
    R_l(r) = A_lr^l + B_lr^{-l-1}
    \]
    定解问题的一般解为：
    \[
    u(r,\theta) = \sum_{l=0}^{\infty}(A_lr^l + B_lr^{-l-1})\mathrm{P}_l(\cos\theta)
    \]
    考虑边界条件：
    \[
    \left.u\right|_{r=0} < \infty,\qquad \left.u\right|_{r=a} = u_0\mathrm{sgn}(\cos\theta)
    \]
    由$\left.u\right|_{r=0} < \infty$得：
    \[
    B_l = 0,\qquad l\in\mathbb{N}
    \]
    将边界条件按本征函数展开：
    \[
    u_0\mathrm{sgn}(\cos\theta) = \sum_{k=0}^{\infty}c_l \mathrm{P}_l(x)
    \]
    由于：
    \[
    \mathrm{sgn}(\cos\theta) = 2\eta(\cos\theta)-1
    \]
    首先考虑将$\eta(\cos\theta)$按Legendre多项式展开：
    \[
    \eta(\cos\theta) = \sum_{k=0}^{\infty}m_l \mathrm{P}_l(x)
    \]
    当$l\neq0$时：
    \begin{align*}
        m_l
        &= \frac{2l+1}{2}\int_{-1}^{1}\eta(\cos\theta)\mathrm{P}_l(\cos\theta)\,\mathrm{d}\cos\theta \\
        &= \frac{2l+1}{2}\int_{0}^{1}\mathrm{P}_l(\cos\theta)\,\mathrm{d}\cos\theta
    \end{align*}
    代入Legendre方程：
    \[
    \frac{\mathrm{d}}{\mathrm{d}\cos\theta}\left[(1-\cos^2\theta)\frac{\mathrm{d}\mathrm{P}_l(\cos\theta)}{\mathrm{d}\cos\theta}\right] + l(l+1)\mathrm{P}_l(\cos\theta) = 0
    \]
    即
    \[
    \mathrm{P}_l(\cos\theta) = \frac{-1}{l(l+1)}\cdot\frac{\mathrm{d}}{\mathrm{d}\cos\theta}\left[(1-\cos^2\theta)\frac{\mathrm{d}\mathrm{P}_l(\cos\theta)}{\mathrm{d}\cos\theta}\right]
    \]
    得：
    \begin{align*}
        m_l
        &= \frac{2l+1}{2}\int_{0}^{1}\mathrm{P}_l(\cos\theta)\,\mathrm{d}\cos\theta \\
        &= \frac{2l+1}{2}\cdot\frac{-1}{l(l+1)}\int_{-1}^{1}\frac{\mathrm{d}}{\mathrm{d}\cos\theta}\left[(1-\cos^2\theta)\frac{\mathrm{d}\mathrm{P}_l(\cos\theta)}{\mathrm{d}\cos\theta}\right]\mathrm{d}\cos\theta \\
        &= \frac{2l+1}{2}\cdot\frac{-1}{l(l+1)}\int_{-1}^{1}\mathrm{d}\left[(1-\cos^2\theta)\frac{\mathrm{d}\mathrm{P}_l(\cos\theta)}{\mathrm{d}\cos\theta}\right] \\
        &= \left.-\frac{2l+1}{2l(l+1)}(1-\cos^2\theta)\frac{\mathrm{d}\mathrm{P}_l(\cos\theta)}{\mathrm{d}\cos\theta}\right|_0^1 \\
        &= \frac{2l+1}{2l(l+1)}\mathrm{P}'_l(0) \\
        &= \left\{
        \begin{matrix}
            &0,\qquad &l = 2k \\
            &\dfrac{4k+3}{2(2k+1)(2k+2)}\cdot\dfrac{(-1)^k (2k+2)!}{2^{2k+1} k!(k+1)!},\qquad &l = 2k+1
        \end{matrix}
        \right.
    \end{align*}
    当$l\neq0$时：
    \[
    \int_{0}^{1}\mathrm{P}_0(\cos\theta)\,\mathrm{d}\cos\theta = 1
    \]
    则有：
    \begin{align*}
    \eta(\cos\theta)
    &= \frac{1}{2} + \frac{1}{2}\sum_{k=0}^{\infty}\frac{4k+3}{(2k+1)(2k+2)}\cdot\frac{(-1)^k (2k+2)!}{2^{2k+1} k!(k+1)!}\mathrm{P}_{2k+1}(\cos\theta) \\
    &= \frac{1}{2} + \frac{1}{2}\sum_{k=0}^{\infty}\frac{4k+3}{2k+2}\cdot\frac{(-1)^k (2k)!}{\left(2^{k}k!\right)^2}\mathrm{P}_{2k+1}(\cos\theta)
    \end{align*}
    于是：
    \[
    \mathrm{sgn}(\cos\theta) = \sum_{k=0}^{\infty}\frac{4k+3}{2k+2}\cdot\frac{(-1)^k (2k)!}{\left(2^{k}k!\right)^2}\mathrm{P}_{2k+1}(\cos\theta)
    \]
    \[
    \left.u\right|_{r=a} = u_0\mathrm{sgn}(\cos\theta) = u_0\sum_{k=0}^{\infty}\frac{4k+3}{2k+2}\cdot\frac{(-1)^k (2k)!}{\left(2^{k}k!\right)^2}\mathrm{P}_{2k+1}(\cos\theta)
    \]
    即：
    \[
    \sum_{l=0}^{\infty}(A_l a^l)\mathrm{P}_l(\cos\theta) = u_0\sum_{k=0}^{\infty}\frac{4k+3}{2k+2}\cdot\frac{(-1)^k (2k)!}{\left(2^{k}k!\right)^2}\mathrm{P}_{2k+1}(\cos\theta)
    \]
    解得：
    \[
    A_l = \left\{
    \begin{aligned}
        &0,\qquad &l = 2k \\
        &\frac{u_0}{a^{2k+1}}\cdot\frac{4k+3}{2k+2}\cdot\frac{(-1)^k (2k)!}{\left(2^{k}k!\right)^2},\qquad &l = 2k+1
    \end{aligned}
    \right.
    \]
    综上，定解问题的解为：
    \[
    u(r,\theta) = u_0\sum_{k=0}^{\infty}\frac{4k+3}{2k+2}\cdot\frac{(-1)^k (2k)!}{\left(2^{k}k!\right)^2}\left(\frac{r}{a}\right)^{2k+1}\mathrm{P}_{2k+1}(\cos\theta)
    \]
    \bigskip

    \item[10.(2)]
    由于只出现$\cos\phi$，说明展开式中一定只含有$m=\pm1$的项，相应地便需要将函数$\sin\theta$与$\cos\theta\sin\theta$设法化为$\mathrm{P}_l^1(\cos\theta)$的形式，作变换$x=\cos\theta$得：
    \[
    \sin\theta = (1-x^2)^{\tfrac{1}{2}}
    \]
    \[
    \cos\theta\sin\theta = x(1-x^2)^{\tfrac{1}{2}}
    \]
    又因为：
    \[
    \mathrm{P}_1^1(x) = \frac{(-1)^1}{2^1 1!}(1-x^2)^{\tfrac{1}{2}}\frac{\mathrm{d}^2}{\mathrm{d}x^2}(x^2-1)^1 = -(1-x^2)^{\tfrac{1}{2}}
    \]
    \[
    \mathrm{P}_2^1(x) = \frac{(-1)^2}{2^2 2!}(1-x^2)^{\tfrac{1}{2}}\frac{\mathrm{d}^3}{\mathrm{d}x^3}(x^2-1)^2 = -3x(1-x^2)^{\tfrac{1}{2}}
    \]
    故：
    \[
    \sin\theta = -\mathrm{P}_1^1(\cos\theta)
    \]
    \[
    \cos\theta\sin\theta = -\frac{1}{3}\mathrm{P}_2^1(\cos\theta)
    \]
    相应的球谐函数$\mathrm{Y}_l^m(\theta,\phi)$为：
    \[
    \mathrm{Y}_1^1(\theta,\phi) = \sqrt{\frac{3}{8\pi}}\mathrm{P}_1^1(\cos\theta)\mathrm{e}^{\mathrm{i}\phi}
    \]
    \[
    \mathrm{Y}_2^1(\theta,\phi) = \sqrt{\frac{5}{24\pi}}\mathrm{P}_2^1(\cos\theta)\mathrm{e}^{\mathrm{i}\phi}
    \]
    则原函数按球谐函数$\mathrm{Y}_l^m(\theta,\phi)$的展开为：
    \begin{align*}
        (1+\cos\theta)\sin\theta\cos\phi
        &= \sin\theta\cos\phi + \cos\theta\sin\theta\cos\phi \\
        &= \frac{1}{2}\sin\theta(\mathrm{e}^{\mathrm{i}\phi} + \mathrm{e}^{-\mathrm{i}\phi}) + \frac{1}{2}\cos\theta\sin\theta(\mathrm{e}^{\mathrm{i}\phi} + \mathrm{e}^{-\mathrm{i}\phi}) \\
        &= -\frac{1}{2}\mathrm{P}_1^1(\cos\theta)(\mathrm{e}^{\mathrm{i}\phi} + \mathrm{e}^{-\mathrm{i}\phi}) - \frac{1}{2}\cdot\frac{1}{3}\mathrm{P}_2^1(\cos\theta)(\mathrm{e}^{\mathrm{i}\phi} + \mathrm{e}^{-\mathrm{i}\phi}) \\
        &= -\frac{1}{2}\sqrt{\frac{8\pi}{3}}\left[\mathrm{Y}_1^1(\theta,\phi) + \mathrm{Y}_1^{-1}(\theta,\phi)\right] - \frac{1}{6}\sqrt{\frac{24\pi}{5}}\left[\mathrm{Y}_2^1(\theta,\phi) + \mathrm{Y}_2^{-1}(\theta,\phi)\right] \\
        &= -\sqrt{\frac{2\pi}{3}}\left[\mathrm{Y}_1^1(\theta,\phi) + \mathrm{Y}_1^{-1}(\theta,\phi)\right] - \sqrt{\frac{2\pi}{15}}\left[\mathrm{Y}_2^1(\theta,\phi) + \mathrm{Y}_2^{-1}(\theta,\phi)\right]
    \end{align*}
    \bigskip

    \item[11.(1)]
    设半球内的温度分布为$u = u(r,\theta,\phi)$，则可得定解问题：
    \[
    \nabla^2 u = \frac{1}{r^2}\frac{\partial}{\partial r}\left(r^2\frac{\partial u}{\partial r}\right) + \frac{1}{r^2\sin\theta}\frac{\partial}{\partial\theta}\left(\sin\theta\frac{\partial u}{\partial\theta}\right) + \frac{1}{r^2\sin^2\theta}\frac{\partial^2 u}{\partial\phi^2} = 0
    \]
    \[
    \left.u\right|_{r=0} < \infty,\qquad \left.u\right|_{r=a} = \mathrm{P}_1^1(\cos\theta)\cos\phi
    \]
    \[
    \left.u\right|_{\theta=0} < \infty,\qquad \left.u\right|_{\theta=\pi} < \infty
    \]
    \[
    \left.u\right|_{\phi=0} = \left.u\right|_{\phi=2\pi},\qquad \left.\frac{\partial u}{\partial\phi}\right|_{\phi=0} = \left.\frac{\partial u}{\partial\phi}\right|_{\phi=2\pi}
    \]
    设方程的解具有分离变量的形式：
    \[
    u(r,\theta) = R(r)S(\theta,\phi)
    \]
    代入方程$\nabla^2 u=0$并分离变量得：
    \[
    \frac{\mathrm{d}}{\mathrm{d}r}\left(r^2\frac{\mathrm{d}R}{\mathrm{d}r}\right) - \lambda R = 0 \tag{1}
    \]
    \[
    \frac{1}{\sin\theta}\cdot\frac{\partial}{\partial\theta}\left(\sin\theta\frac{\partial S(\theta,\phi)}{\partial\theta}\right) + \frac{1}{\sin^2\theta}\cdot\frac{\partial^2 S(\theta,\phi)}{\partial\phi^2} + \lambda S(\theta,\phi) = 0 \tag{2}
    \]
    由$\theta$和$\phi$的边界条件，得：
    \[
    \left.S\right|_{\theta=0} < \infty,\qquad \left.S\right|_{\theta=\pi} < \infty \tag{3}
    \]
    \[
    \left.S\right|_{\phi=0} = \left.S\right|_{\phi=2\pi},\qquad \left.\frac{\partial S}{\partial\phi}\right|_{\phi=0} = \left.\frac{\partial S}{\partial\phi}\right|_{\phi=2\pi} \tag{4}
    \]
    (2)、(3)和(4)构成了本征值问题，其解为：
    \[
    \lambda_l = l(l+1)
    \]
    \[
    S_l^m(\theta,\phi) = \mathrm{Y}_l^m(\theta,\phi),\qquad m = 0,\pm1,\pm2,...,\pm l
    \]
    代入方程(1)得：
    \[
    \frac{\mathrm{d}}{\mathrm{d}r}\left(r^2\frac{\mathrm{d}R}{\mathrm{d}r}\right) - l(l+1)R = 0 \tag{5}
    \]
    其解为：
    \[
    R_l(r) = A_l^m r^l + B_l^m r^{-l-1}
    \]
    定解问题的一般解为：
    \[
    u(r,\theta) = \sum_{l=0}^{\infty}\sum_{m=-l}^{l}\left(A_l^m r^l + B_l^m r^{-l-1}\right)\mathrm{Y}_l^m(\theta,\phi)
    \]
    考虑边界条件：
    \[
    \left.u\right|_{r=0} < \infty,\qquad \left.u\right|_{r=a} = \mathrm{P}_1^1(\cos\theta)\cos\phi
    \]
    由$\left.u\right|_{r=0} < \infty$得：
    \[
    B_l = 0,\qquad l\in\mathbb{N}
    \]
    将边界条件按球谐函数展开：
    \[
    \mathrm{P}_1^1(\cos\theta)\cos\phi = \sqrt{\frac{2\pi}{3}}\left[\mathrm{Y}_1^1(\theta,\phi) + \mathrm{Y}_1^{-1}(\theta,\phi)\right]
    \]
    于是：
    \[
    \left.u\right|_{r=a} = \mathrm{P}_1^1(\cos\theta)\cos\phi = \sqrt{\frac{2\pi}{3}}\left[\mathrm{Y}_1^1(\theta,\phi) + \mathrm{Y}_1^{-1}(\theta,\phi)\right]
    \]
    即：
    \[
    \sum_{l=0}^{\infty}\sum_{m=-l}^{l}A_l^m a^l \mathrm{Y}_l^m(\theta,\phi) = \sqrt{\frac{2\pi}{3}}\left[\mathrm{Y}_1^1(\theta,\phi) + \mathrm{Y}_1^{-1}(\theta,\phi)\right]
    \]
    解得：
    \[
    A_1^1 = A_1^{-1} = \frac{1}{a}\sqrt{\frac{2\pi}{3}}
    \]
    其余系数均为0。

    综上，定解问题的解为：
    \begin{align*}
    u(r,\theta)
    &= \frac{r}{a}\sqrt{\frac{2\pi}{3}}\left[\mathrm{Y}_1^1(\theta,\phi) + \mathrm{Y}_1^{-1}(\theta,\phi)\right] \\
    &= \frac{r}{a}\mathrm{P}_1^1(\cos\theta)\cos\phi \\
    &= -\frac{r}{a}\sin\theta\cos\phi
    \end{align*}
    \bigskip

    \item[11.(2)]
    设半球内的温度分布为$u = u(r,\theta,\phi)$，则可得定解问题：
    \[
    \nabla^2 u = \frac{1}{r^2}\frac{\partial}{\partial r}\left(r^2\frac{\partial u}{\partial r}\right) + \frac{1}{r^2\sin\theta}\frac{\partial}{\partial\theta}\left(\sin\theta\frac{\partial u}{\partial\theta}\right) + \frac{1}{r^2\sin^2\theta}\frac{\partial^2 u}{\partial\phi^2} = 0
    \]
    \[
    \left.u\right|_{r=0} < \infty,\qquad \left.u\right|_{r=a} = \mathrm{P}_1(\cos\theta)\sin\theta\cos\phi
    \]
    \[
    \left.u\right|_{\theta=0} < \infty,\qquad \left.u\right|_{\theta=\pi} < \infty
    \]
    \[
    \left.u\right|_{\phi=0} = \left.u\right|_{\phi=2\pi},\qquad \left.\frac{\partial u}{\partial\phi}\right|_{\phi=0} = \left.\frac{\partial u}{\partial\phi}\right|_{\phi=2\pi}
    \]
    设方程的解具有分离变量的形式：
    \[
    u(r,\theta) = R(r)S(\theta,\phi)
    \]
    代入方程$\nabla^2 u=0$并分离变量得：
    \[
    \frac{\mathrm{d}}{\mathrm{d}r}\left(r^2\frac{\mathrm{d}R}{\mathrm{d}r}\right) - \lambda R = 0 \tag{1}
    \]
    \[
    \frac{1}{\sin\theta}\cdot\frac{\partial}{\partial\theta}\left(\sin\theta\frac{\partial S(\theta,\phi)}{\partial\theta}\right) + \frac{1}{\sin^2\theta}\cdot\frac{\partial^2 S(\theta,\phi)}{\partial\phi^2} + \lambda S(\theta,\phi) = 0 \tag{2}
    \]
    由$\theta$和$\phi$的边界条件，得：
    \[
    \left.S\right|_{\theta=0} < \infty,\qquad \left.S\right|_{\theta=\pi} < \infty \tag{3}
    \]
    \[
    \left.S\right|_{\phi=0} = \left.S\right|_{\phi=2\pi},\qquad \left.\frac{\partial S}{\partial\phi}\right|_{\phi=0} = \left.\frac{\partial S}{\partial\phi}\right|_{\phi=2\pi} \tag{4}
    \]
    (2)、(3)和(4)构成了本征值问题，其解为：
    \[
    \lambda_l = l(l+1)
    \]
    \[
    S_l^m(\theta,\phi) = \mathrm{Y}_l^m(\theta,\phi),\qquad m = 0,\pm1,\pm2,...,\pm l
    \]
    代入方程(1)得：
    \[
    \frac{\mathrm{d}}{\mathrm{d}r}\left(r^2\frac{\mathrm{d}R}{\mathrm{d}r}\right) - l(l+1)R = 0 \tag{5}
    \]
    其解为：
    \[
    R_l(r) = A_l^m r^l + B_l^m r^{-l-1}
    \]
    定解问题的一般解为：
    \[
    u(r,\theta) = \sum_{l=0}^{\infty}\sum_{m=-l}^{l}\left(A_l^m r^l + B_l^m r^{-l-1}\right)\mathrm{Y}_l^m(\theta,\phi)
    \]
    考虑边界条件：
    \[
    \left.u\right|_{r=0} < \infty,\qquad \left.u\right|_{r=a} = \mathrm{P}_1(\cos\theta)\sin\theta\cos\phi
    \]
    由$\left.u\right|_{r=0} < \infty$得：
    \[
    B_l = 0,\qquad l\in\mathbb{N}
    \]
    将边界条件按球谐函数展开：
    \[
    \mathrm{P}_1(\cos\theta)\sin\theta\cos\phi = \cos\theta\sin\theta\cos\phi = -\sqrt{\frac{2\pi}{15}}\left[\mathrm{Y}_2^1(\theta,\phi) + \mathrm{Y}_2^{-1}(\theta,\phi)\right]
    \]
    于是：
    \[
    \left.u\right|_{r=a} = \mathrm{P}_1(\cos\theta)\sin\theta\cos\phi = -\sqrt{\frac{2\pi}{15}}\left[\mathrm{Y}_2^1(\theta,\phi) + \mathrm{Y}_2^{-1}(\theta,\phi)\right]
    \]
    即：
    \[
    \sum_{l=0}^{\infty}\sum_{m=-l}^{l}A_l^m a^l \mathrm{Y}_l^m(\theta,\phi) = -\sqrt{\frac{2\pi}{15}}\left[\mathrm{Y}_2^1(\theta,\phi) + \mathrm{Y}_2^{-1}(\theta,\phi)\right]
    \]
    解得：
    \[
    A_2^1 = A_2^{-1} = -\frac{1}{a}\sqrt{\frac{2\pi}{15}}
    \]
    其余系数均为0。

    综上，定解问题的解为：
    \begin{align*}
        u(r,\theta)
        &= -\left(\frac{r}{a}\right)^2\sqrt{\frac{2\pi}{15}}\left[\mathrm{Y}_2^1(\theta,\phi) + \mathrm{Y}_2^{-1}(\theta,\phi)\right] \\
        &= -\frac{1}{3}\left(\frac{r}{a}\right)^2 \mathrm{P}_2^1(\cos\theta)\cos\phi \\
        &= \left(\frac{r}{a}\right)^2\cos\theta\sin\theta\cos\phi
    \end{align*}
    \bigskip

    \item[12.]
    设球内问题的解可以写为如下形式：
    \[
    u(r,\theta,\phi) = v(r,\theta,\phi) + w(r,\theta,\phi)
    \]
    其中：
    \[
    \nabla^2 v = A
    \]
    \[
    \left.u\right|_{r=0} < \infty,\qquad \left.u\right|_{r=a} = 0
    \]
    且
    \[
    \nabla^2 w = Br^2\sin{2\theta}\cos\phi
    \]
    \[
    \left.w\right|_{r=0} < \infty,\qquad \left.w\right|_{r=a} = 0
    \]
    下面考虑这样的定解问题：
    \[
    \nabla^2 U = Kr^n \mathrm{P}_l^m(cos\theta)\cos{m\phi}
    \]
    \[
    \left.U\right|_{r=0} < \infty,\qquad \left.U\right|_{r=a} = 0
    \]
    \[
    \left.U\right|_{\theta=0} < \infty,\qquad \left.U\right|_{\theta=\pi} < \infty
    \]
    \[
    \left.U\right|_{\phi=0} = \left.U\right|_{\phi=2\pi},\qquad \left.\frac{\partial U}{\partial\phi}\right|_{\phi=0} = \left.\frac{\partial U}{\partial\phi}\right|_{\phi=2\pi}
    \]
    考虑到$\theta$和$\phi$的边界条件，设方程的解有如下形式：
    \[
    U(r,\theta,\phi) = \sum_{k=0}^{\infty}\sum_{s=0}^{k}\left[M_k^s(r)\sin{s\phi} + N_k^s(r)\cos{s\phi}\right]\mathrm{P}_s^k(\cos\theta)
    \]
    代入$U$所满足的方程及$r$的边界条件，并利用本征函数的正交性，即可推得$M_k^s(r)$和$N_k^s(r)$满足的定解问题：
    \[
    \frac{\mathrm{d}}{\mathrm{d}r}\left(r^2\frac{\mathrm{d}M_k^s}{\mathrm{d}r}\right) - k(k+1)M_k^s = 0 \tag{1}
    \]
    \[
    \left.M_k^s\right|_{r=0} < \infty,\qquad \left.M_k^s\right|_{r=a} = 0 \tag{2}
    \]
    以及
    \[
    \frac{\mathrm{d}}{\mathrm{d}r}\left(r^2\frac{\mathrm{d}N_k^s}{\mathrm{d}r}\right) - k(k+1)N_k^s = Ar^{n+2}\delta_{kl}\delta_{sm} \tag{3}
    \]
    \[
    \left.N_k^s\right|_{r=0} < \infty,\qquad \left.N_k^s\right|_{r=a} = 0 \tag{4}
    \]
    方程(1)的通解为：
    \[
    M_k^s(r) = Cr^k + Dr^{-k-1}
    \]
    代入边界条件(2)后，可以得到：
    \[
    C = D = 0
    \]
    即：
    \[
    M_k^s(r) = 0
    \]
    显然，对于方程(3)和边界条件(4)，当$k \neq l$、$s \neq m$时，定解问题退化为(1)和(2)；当$k=l$、$s=m$时，方程才有非零解。

    分两种情况讨论：

    $n \neq l-2$时，方程(3)的解不含对数项：
    \[
    N_k^s = \frac{K}{(n+2)(n+3)-l(l+1)}r^{n+2} + Er^l + Fr^{-l-1}
    \]
    代入边界条件(4)，可得：
    \[
    E = -\frac{K}{(n+2)(n+3)-l(l+1)}a^{n-l+2},\qquad F = 0
    \]
    \[
    N_k^s = \frac{K}{(n+2)(n+3)-l(l+1)}\left(r^{n+2} - a^{n-l+2}r^l\right)
    \]
    $n = l-2$时，方程(3)的解必含对数项：
    \[
    N_k^s = \frac{K}{2l+1}r^l\ln{r} + G
    \]
    代入边界条件(4)，可得：
    \[
    G = -\frac{K}{2l+1}r^l\ln{a}
    \]
    \[
    N_k^s = \frac{K}{2l+1}r^l(\ln{r} - \ln{a})
    \]
    故关于$U$的定解问题的解为：
    \[
    U(r,\theta,\phi) = \left\{
    \begin{aligned}
        &\frac{K}{(n+2)(n+3)-l(l+1)}\left(r^{n+2} - a^{n-l+2}r^l\right)\mathrm{P}_l^m(\cos\theta)\cos{m\phi},\qquad &n \neq l-2 \\
        &\frac{K}{2l+1}r^l(\ln{r} - \ln{a})\mathrm{P}_l^m(\cos\theta)\cos{m\phi},\qquad &n = l-2
    \end{aligned}
    \right.
    \]
    对于$v(r,\theta,\phi)$，$K = A$，$n = m = l = 0$，且$n \neq l-2$，定解问题的解为：
    \[
    v(r,\theta,\phi) = \frac{A}{6}(r^2 - a^2)
    \]
    对于$w(r,\theta,\phi)$，注意到：
    \begin{align*}
    \nabla^2 w
    &= Br^2\sin{2\theta}\cos\phi \\
    &= 2Br^2\cos\theta\sin\theta\cos\phi \\
    &= -\frac{2}{3}Br^2\mathrm{P}_2^1(\cos\theta)\cos\phi
    \end{align*}
    则$K = -\dfrac{2}{3}B$，$n = l = 2$，$m = 1$，且$n \neq l-2$，定解问题的解为：
    \[
    w(r,\theta,\phi) = -\frac{B}{21}r^2(r^2 - a^2)\mathrm{P}_2^1(\cos\theta)\cos\phi
    \]
    综上，原球内问题的解为：
    \begin{align*}
        u(r,\theta,\phi)
        &= v(r,\theta,\phi) + w(r,\theta,\phi) \\
        &= \frac{A}{6}(r^2 - a^2) - \frac{B}{21}r^2(r^2 - a^2)\mathrm{P}_2^1(\cos\theta)\cos\phi \\
        &= \frac{A}{6}(r^2 - a^2) + \frac{B}{14}r^2(r^2 - a^2)\sin{2\theta}\cos\phi
    \end{align*}

    \end{itemize}

\end{document}
